% ============================================================
% Section 114: 自旋1量子气体与经典偶极子的顺磁性
% ============================================================
\section{固体物理：自旋为1的量子气体与经典偶极子顺磁性}
\label{sec:spin1-classical-paramagnet}

\begin{modelbox}[题目：量子与经典顺磁磁化率]
考虑一稳定的矢性介子（自旋为1的重粒子）的理想气体，粒子磁矩为 $\mu$。

(1) 计算弱场下单位体积中的顺磁磁化率；

(2) 对经典偶极子作相应的计算。
\end{modelbox}

\begin{tipbox}[考点快速分析]
\begin{itemize}
    \item \textbf{自旋 $S=1$}：$m_S=-1,0,+1$ 三个投影，配分函数 $Z_1=2\cosh(\beta\mu B)+1$
    \item \textbf{经典偶极子}：连续取向，朗之万函数 $L(y)=\coth y-1/y$
    \item \textbf{弱场极限}：两者均给出居里定律 $\chi\propto1/T$
\end{itemize}
\end{tipbox}

\begin{derivationbox}[标准解题过程]
\textbf{(1) 自旋为1的量子气体的顺磁磁化率}

自旋 $S=1$ 的粒子在磁场 $B$ 中，磁矩沿磁场方向的投影有3个可能取值：$+\mu$、$0$、$-\mu$。对应的附加能量为 $\epsilon=-\mu B,0,+\mu B$。

单个粒子的配分函数：
\begin{equation}
Z_1=e^{\beta\mu B}+1+e^{-\beta\mu B}=2\cosh(\beta\mu B)+1,
\end{equation}
其中 $\beta=1/(k_BT)$。

沿磁场方向的平均磁矩：
\begin{equation}
\bar{\mu}=\frac{\mu e^{\beta\mu B}+0\cdot1-\mu e^{-\beta\mu B}}{Z_1}
=\mu\,\frac{2\sinh(\beta\mu B)}{2\cosh(\beta\mu B)+1}.
\end{equation}

设单位体积粒子数为 $n$，则磁化强度 $M=n\bar{\mu}$。

弱场极限 $\mu B\ll k_BT$（即 $\beta\mu B\ll1$）：
\begin{equation}
\sinh(\beta\mu B)\approx\beta\mu B,\qquad \cosh(\beta\mu B)\approx1.
\end{equation}
因此
\begin{equation}
\bar{\mu}\approx\mu\cdot\frac{2\beta\mu B}{3}=\frac{2\mu^2B}{3k_BT},
\end{equation}
\begin{equation}
M\approx\frac{2n\mu^2}{3k_BT}B.
\end{equation}

顺磁磁化率（SI 单位制 $\chi=\mu_0M/B$；高斯制 $\chi=M/B$）：
\begin{equation}
\boxed{\chi=\frac{2n\mu_0\mu^2}{3k_BT}\quad(\text{SI}),\qquad\text{或}\qquad\chi=\frac{2n\mu^2}{3k_BT}\quad(\text{CGS}).}
\end{equation}
此即自旋为1理想顺磁气体的\textbf{居里定律}，居里常数 $C=2n\mu_0\mu^2/(3k_B)$。

\textbf{(2) 经典偶极子的顺磁磁化率}

经典偶极子的磁矩 $\bm{\mu}$ 可连续取向，与外磁场 $B$ 夹角为 $\theta$。附加能量
\begin{equation}
\epsilon=-\bm{\mu}\cdot\bm{B}=-\mu B\cos\theta.
\end{equation}

沿磁场方向的平均磁矩由玻尔兹曼分布加权平均：
\begin{equation}
\bar{\mu}=\frac{\int_0^\pi\mu\cos\theta\,e^{\beta\mu B\cos\theta}\,2\pi\sin\theta\,d\theta}{\int_0^\pi e^{\beta\mu B\cos\theta}\,2\pi\sin\theta\,d\theta}.
\end{equation}

令 $x=\cos\theta$，$y=\beta\mu B$，积分限 $-1\le x\le1$：
\begin{equation}
\bar{\mu}=\mu\,\frac{\int_{-1}^1x\,e^{yx}\,dx}{\int_{-1}^1e^{yx}\,dx}
=\mu\left(\coth y-\frac{1}{y}\right)\equiv\mu L(y),
\end{equation}
其中 $L(y)=\coth y-1/y$ 为\textbf{朗之万函数}。

弱场极限 $y=\mu B/k_BT\ll1$：
\begin{equation}
\coth y=\frac{1}{y}+\frac{y}{3}-\frac{y^3}{45}+\cdots,
\end{equation}
故 $L(y)\approx y/3$。因此
\begin{equation}
\bar{\mu}\approx\mu\cdot\frac{\mu B}{3k_BT}=\frac{\mu^2B}{3k_BT},
\end{equation}
\begin{equation}
M=n\bar{\mu}=\frac{n\mu^2}{3k_BT}B.
\end{equation}

顺磁磁化率：
\begin{equation}
\boxed{\chi=\frac{n\mu_0\mu^2}{3k_BT}\quad(\text{SI}),\qquad\text{或}\qquad\chi=\frac{n\mu^2}{3k_BT}\quad(\text{CGS}).}
\end{equation}
此即经典朗之万顺磁性的居里定律。
\end{derivationbox}

\begin{formulabox}[答案汇总]
\begin{center}
\begin{tabular}{lll}
\toprule
\textbf{系统} & \textbf{弱场磁化率 $\chi$（SI）} & \textbf{居里常数} \\
\midrule
自旋 $S=1$ 量子气体 & $\displaystyle\frac{2n\mu_0\mu^2}{3k_BT}$ & $\displaystyle\frac{2n\mu_0\mu^2}{3k_B}$ \\
经典偶极子气体 & $\displaystyle\frac{n\mu_0\mu^2}{3k_BT}$ & $\displaystyle\frac{n\mu_0\mu^2}{3k_B}$ \\
\bottomrule
\end{tabular}
\end{center}
\end{formulabox}

\begin{insightbox}[物理图像]
\begin{itemize}
    \item 自旋 $S=1$ 系统的磁化率是经典偶极子的\textbf{2 倍}。这是因为量子系统中存在一个零投影态（$m_S=0$），使得磁矩在磁场方向的有效涨落比经典连续取向更大。
    \item 对于有限自旋 $S$，弱场极限下 $\chi\propto S(S+1)$。自旋 $1/2$ 时 $\chi\propto1/4$，自旋 $1$ 时 $\chi\propto2/3$（以统一的磁矩标度 $\mu$ 衡量）。
    \item 两种系统在弱场均遵从居里定律 $\chi\propto1/T$，反映热运动对磁矩取向的扰乱。经典极限对应 $S\to\infty$ 的布里渊函数行为。
\end{itemize}
\end{insightbox}
